\documentclass[10pt,a4paper]{article}
\usepackage[margin=25mm]{geometry}
\usepackage{amsmath,amssymb,lmodern}
\setlength{\emergencystretch}{2em}
\begin{document}
\begin{center}
{\Large A Weyl extremal equality without a common eigenbasis}\\[4pt]
{\large Counterexample to Conjecture 00000007152}\\[6pt]
ziangni-sys \qquad 4 October 2026
\end{center}

\section*{The extremal equality}
For real symmetric matrices, the largest-eigenvalue Weyl inequality is
\[
\lambda_{\max}(A+B)\leq\lambda_{\max}(A)+\lambda_{\max}(B).
\]
We refute the source's necessity of a common eigenbasis for Weyl extremal
equality by exhibiting equality in this particular bound without such
a basis. The example does not assert simultaneous saturation of every
Weyl inequality; that is a stronger condition than equality in this
extremal bound.

Consider the actual symmetric matrices and vector
\[
A=\begin{pmatrix}3&0&0\\0&1&0\\0&0&0\end{pmatrix},\qquad
B=\begin{pmatrix}3&0&0\\0&1&1\\0&1&1\end{pmatrix},\qquad
e=\begin{pmatrix}1\\0\\0\end{pmatrix}.
\]
Both $Ae=3e$ and $Be=3e$, so $(A+B)e=6e$.

\section*{Actual largest eigenvalues}
Write $v=(v_0,v_1,v_2)$ and $S(v)=v_0^2+v_1^2+v_2^2$.
Then
\[
v^TAv=3v_0^2+v_1^2\leq3S(v),\qquad
v^TBv=3v_0^2+(v_1+v_2)^2\leq3S(v).
\]
The second bound follows from
$(v_1+v_2)^2\leq2(v_1^2+v_2^2)$, since $(v_1-v_2)^2\geq0$.
Adding the two bounds gives $v^T(A+B)v\leq6S(v)$.
For any nonzero real eigenvector $Mv=\mu v$, the quadratic form is
$v^TMv=\mu S(v)$ and $S(v)>0$. These bounds therefore imply that
every eigenvalue of $A$, $B$, and $A+B$ is at most $3$, $3$, and $6$,
respectively. Together with the displayed eigenvector $e$, they prove
\[
\lambda_{\max}(A)=3,\qquad\lambda_{\max}(B)=3,\qquad
\lambda_{\max}(A+B)=6=3+3.
\]
This verifies the required equality using the matrices themselves,
rather than assuming an eigenvalue table.

\section*{No common eigenbasis}
If a basis $\{b_i\}$ consisted of simultaneous eigenvectors, with
$Ab_i=\alpha_i b_i$ and $Bb_i=\beta_i b_i$, then
\[
ABb_i=\beta_i\alpha_i b_i=\alpha_i\beta_i b_i=BAb_i.
\]
Linearity and the spanning property of the basis would give $AB=BA$.
However, for $w=(0,0,1)^T$ direct computation gives
\[
ABw=(0,1,0)^T,\qquad BAw=(0,0,0)^T.
\]
Thus no common eigenbasis exists, even without requiring the basis to
be orthogonal. The matrices share the one extremal eigenvector needed
to saturate the top Weyl bound, while their actions on its complementary
plane do not commute. A common eigenbasis is consequently not necessary
for this Weyl extremal equality.

\section*{Formal verification}
Lean defines actual matrix multiplication, eigenvalue witnesses, and
universal largest-eigenvalue bounds. It checks symmetry and all three
largest eigenvalues. The theorem \texttt{no\_common\_eigenbasis} quantifies
over an actual Lean \texttt{Basis} with arbitrary index type and proves
the commutation contradiction. The final theorem packages the exact
extremal equality together with the absence of a common eigenbasis.
\end{document}
